% =============================================================
%  MAT231-F27 Calculus II — Lecture 05
%  Unit 5: Power series; radius of convergence; term-by-term
%  differentiation and integration
%
%  Build from the repository root:   npm run materials -- lecture-05
%  This writes lecture-05.pdf (slides, with overlays) and
%  lecture-05-handout.pdf (print version) into this folder.
% =============================================================
\documentclass{wosm-slides}

\newcommand{\R}{\mathbb{R}}
\newcommand{\dt}{\,\mathrm{d}t}

\course{MAT231-F27}{Calculus II}
\title[L05 \textperiodcentered\ Power series]{Power series}
\subtitle{Radius of convergence; differentiation and integration term by term}
\author{Dr.\ Ali Khudiyev}
\institute{WoSM Academy}
\date{Lecture 05\enspace\textperiodcentered\enspace Unit 5\enspace\textperiodcentered\enspace Fall 2027}

\begin{document}

\maketitle

% -------------------------------------------------------------
\begin{frame}{Last time}
  \begin{itemize}
    \item $f_n \to f$ \textbf{uniformly} on $D$ if $\sup_{x \in D} |f_n(x) - f(x)| \to 0$.
    \item Uniform limits of continuous functions are continuous; on $[a, b]$ we may swap
      $\lim$ and $\int$; derivatives need uniform convergence of $f_n'$.
    \item \textbf{M-test}: $|f_k(x)| \le M_k$ with $\sum M_k < \infty$ gives uniform convergence
      of $\sum f_k$.
    \item $\sum x^k = \frac{1}{1-x}$ converges uniformly on every $[-a, a]$ with $a < 1$, so we
      could integrate and differentiate it term by term.
  \end{itemize}
  \vskip0.6em
  Today: the same for every power series.
\end{frame}

\begin{frame}{This lecture}
  \tableofcontents
  \vskip1em
  {\small\color{wosm-muted}
  Reading: G.~Strang, E.~Herman, \emph{Calculus, Volume 2} (OpenStax), Sections 6.1 and 6.2.\par}
\end{frame}

% =============================================================
\section{Power series}

\begin{frame}{Power series}
  \begin{definition}
    A \textbf{power series centred at $a$} is a series
    \[ \sum_{k=0}^\infty c_k (x - a)^k = c_0 + c_1 (x - a) + c_2 (x - a)^2 + \cdots \]
    with \textbf{coefficients} $c_k \in \R$. (We read $(x - a)^0$ as $1$, also at $x = a$.)
  \end{definition}
  \begin{itemize}
    \item A polynomial ``of infinite degree''. Its partial sums are polynomials.
    \item It always converges at $x = a$, to $c_0$.
    \item For which $x$ does it converge? What function does it define?
      Can we differentiate and integrate it like a polynomial?
  \end{itemize}
\end{frame}

\begin{frame}{Four examples}
  With the ratio test (Lecture 3), for $x \ne 0$:
  \begin{itemize}
    \item $\displaystyle\sum x^k$:
      \uncover<2->{ratio $|x|$. Converges for $|x| < 1$ only.}
    \item $\displaystyle\sum \frac{x^k}{k!}$:
      \uncover<3->{ratio $\frac{|x|}{k+1} \to 0$. Converges for every $x$.}
    \item $\displaystyle\sum k!\, x^k$:
      \uncover<4->{ratio $(k+1)|x| \to \infty$. Converges only at $x = 0$.}
    \item $\displaystyle\sum_{k=1}^\infty \frac{x^k}{k}$:
      \uncover<5->{ratio $\frac{k}{k+1}|x| \to |x|$. Converges for $|x| < 1$ and diverges for
      $|x| > 1$. At $x = 1$: the harmonic series, divergent. At $x = -1$: the alternating
      harmonic series, convergent. So: $[-1, 1)$.}
  \end{itemize}
  \uncover<6->{In each case the set of convergence is an interval centred at $0$
  --- possibly a single point or all of $\R$.}
\end{frame}

% =============================================================
\section{The radius of convergence}

\begin{frame}{Convergence spreads inwards}
  \begin{lemma}[centre $a = 0$; otherwise replace $x$ by $x - a$]
    If $\sum c_k x_0^k$ converges for some $x_0 \ne 0$, then $\sum c_k x^k$ converges absolutely
    for every $|x| < |x_0|$, and uniformly on $[-r, r]$ for every $r < |x_0|$.
  \end{lemma}
  \pause
  \begin{proof}
    The terms of a convergent series tend to $0$, so they are bounded:
    $|c_k x_0^k| \le M$ for all $k$. For $|x| \le r$,
    $|c_k x^k| = |c_k x_0^k| \, \bigl|\frac{x}{x_0}\bigr|^k \le M q^k$ with
    $q = \frac{r}{|x_0|} < 1$. As $\sum M q^k$ converges, the M-test applies.
  \end{proof}
\end{frame}

\begin{frame}{The radius of convergence}
  \begin{theorem}
    For every power series $\sum c_k (x - a)^k$ there is an $R$ with $0 \le R \le \infty$,
    the \textbf{radius of convergence}, such that the series
    \begin{itemize}
      \item converges absolutely if $|x - a| < R$, uniformly on every $[a - r, a + r]$ with $r < R$;
      \item diverges if $|x - a| > R$.
    \end{itemize}
  \end{theorem}
  {\small Proof: let $R$ be the supremum of $|x - a|$ over all $x$ where the series converges,
  and apply the lemma.\par}
  \pause
  {\centering
  \begin{tikzpicture}[x=1cm, y=1cm]
    \draw[wosm-muted] (-5.5,0) -- (5.5,0);
    \draw[wosm-accent, line width=2.4pt] (-2.6,0) -- (2.6,0);
    \foreach \p in {-2.6, 2.6}
      \node[circle, draw=wosm-text, fill=wosm-bg, inner sep=1.6pt] at (\p,0) {};
    \fill[wosm-text] (0,0) circle (1.4pt);
    \node[below, font=\small] at (-2.6,-0.1) {$a - R$};
    \node[below, font=\small] at (0,-0.1) {$a$};
    \node[below, font=\small] at (2.6,-0.1) {$a + R$};
    \node[above, font=\small, text=wosm-accent] at (0,0.12) {converges absolutely};
    \node[above, font=\small, text=wosm-muted] at (-4.1,0.12) {diverges};
    \node[above, font=\small, text=wosm-muted] at (4.1,0.12) {diverges};
    \node[above, font=\small, text=wosm-brick] at (-2.6,0.12) {?};
    \node[above, font=\small, text=wosm-brick] at (2.6,0.12) {?};
  \end{tikzpicture}\par}
\end{frame}

\begin{frame}{Computing $R$}
  \begin{theorem}
    If $c_k \ne 0$ and $\Bigl|\dfrac{c_{k+1}}{c_k}\Bigr| \to L$, then $R = \dfrac1L$
    \ (with $\frac10 = \infty$ and $\frac1\infty = 0$).
    The same holds if $\sqrt[k]{|c_k|} \to L$.
  \end{theorem}
  {\small Proof: the ratio of consecutive terms is
  $\Bigl|\frac{c_{k+1}(x-a)^{k+1}}{c_k (x-a)^k}\Bigr| \to L\,|x - a|$, which is below $1$ if and
  only if $|x - a| < \frac1L$. Apply the ratio test (Lecture 3).\par}
  \pause
  \begin{remark}
    The \textbf{Cauchy--Hadamard formula} $R = 1 / \limsup_k \sqrt[k]{|c_k|}$ always holds,
    even when the limits above do not exist. We will not need it.
  \end{remark}
\end{frame}

\begin{frame}{Examples}
  \begin{itemize}
    \item $\displaystyle\sum_{k=1}^\infty \frac{(x-2)^k}{k\, 3^k}$
      \uncover<2->{\\ $\bigl|\frac{c_{k+1}}{c_k}\bigr| = \frac{k}{3(k+1)} \to \frac13$, so $R = 3$:
      converges for $-1 < x < 5$. At $x = 5$: $\sum \frac1k$, divergent; at $x = -1$:
      $\sum \frac{(-1)^k}{k}$, convergent. Interval: $[-1, 5)$.}
    \item $\displaystyle\sum_{k=0}^\infty \frac{x^{2k}}{4^k}$
      \uncover<3->{\\ Every other coefficient is $0$, so the formula does not apply. Ratio test
      on the terms: $\frac{x^2}{4} < 1$ iff $|x| < 2$, so $R = 2$. (Sum: $\frac{1}{1 - x^2/4}$.)}
    \item $\displaystyle\sum_{k=1}^\infty k^k x^k$
      \uncover<4->{\\ $\sqrt[k]{|c_k|} = k \to \infty$, so $R = 0$: it converges only at $x = 0$.}
  \end{itemize}
\end{frame}

\begin{frame}{At the endpoints anything can happen}
  \begin{center}
  \small\renewcommand{\arraystretch}{1.2}
  \begin{tabular}{@{}llll@{}}
    \toprule
    \textbf{Series} & $R$ & \textbf{At $x = 1$} & \textbf{At $x = -1$} \\
    \midrule
    $\sum x^k$               & $1$ & diverges & diverges \\
    $\sum \frac{x^k}{k}$     & $1$ & diverges & converges conditionally \\
    $\sum \frac{(-1)^k x^k}{k}$ & $1$ & converges conditionally & diverges \\
    $\sum \frac{x^k}{k^2}$   & $1$ & converges absolutely & converges absolutely \\
    \bottomrule
  \end{tabular}
  \end{center}
  \vskip0.3em
  Check the endpoints separately, with the tests of Lectures 2 and 3.
  \pause
  \begin{remark}[Abel's theorem, without proof]
    If the series converges at $x = a + R$, its sum is continuous there from the left:
    $\lim_{x \to (a+R)^-} \sum c_k (x-a)^k = \sum c_k R^k$. Likewise at $a - R$.
  \end{remark}
\end{frame}

% =============================================================
\section{Calculus with power series}

\begin{frame}{Differentiation and integration term by term}
  \begin{theorem}
    Let $\sum c_k (x - a)^k$ have radius $R > 0$ and sum $f(x)$ for $|x - a| < R$. Then
    $f$ is continuous and differentiable there, and
    \begin{align*}
      f'(x) &= \sum_{k=1}^\infty k\, c_k (x - a)^{k-1}, &
      \int_a^x f(t)\dt &= \sum_{k=0}^\infty \frac{c_k}{k+1} (x - a)^{k+1} .
    \end{align*}
    Both new series have the same radius $R$.
  \end{theorem}
  \pause
  \begin{itemize}
    \item Inside the interval of convergence a power series behaves like a polynomial.
    \item The radius stays the same; the behaviour at the endpoints may change.
  \end{itemize}
\end{frame}

\begin{frame}{Proof}
  \begin{proof}
    Take $a = 0$ and $|x| < r < R$. The series converges at $r$, so $|c_k r^k| \le M$. Then
    \[ |k c_k x^{k-1}| \le \frac{M}{r}\, k\, q^{k-1}, \qquad q = \frac{|x|}{r} < 1 , \]
    and $\sum k q^{k-1}$ converges (ratio test).
    So the derivative series converges absolutely for $|x| < R$, and by the M-test uniformly on
    every $[-r', r']$ with $r' < R$. Lecture 4: we may differentiate term by term there.
    For integrals: on $[0, x]$ the series converges uniformly; swap $\int$ and $\sum$.
    The radius cannot grow: $|c_k x^k| \le |x| \cdot |k c_k x^{k-1}|$ for $k \ge 1$.
  \end{proof}
\end{frame}

\begin{frame}{Coefficients are derivatives}
  \begin{corollary}
    If $f(x) = \sum c_k (x - a)^k$ for $|x - a| < R$ with $R > 0$, then $f$ has derivatives of
    every order there, and
    \[ c_k = \frac{f^{(k)}(a)}{k!} \qquad (k = 0, 1, 2, \dots) . \]
  \end{corollary}
  \pause
  {\small Proof: differentiate $k$ times. All terms with $j < k$ disappear; the term with $j = k$
  becomes $k!\, c_k$; the terms with $j > k$ still contain a factor $(x - a)$, which is $0$ at
  $x = a$.\par}
  \pause
  \begin{itemize}
    \item \textbf{Uniqueness}: two power series centred at $a$ that agree near $a$ have the same
      coefficients. However we find a power series for $f$, it is the same one.
    \item Next lecture: start from $f$ and use $c_k = \frac{f^{(k)}(a)}{k!}$ --- the Taylor series.
  \end{itemize}
\end{frame}

\begin{frame}{Endpoints can change}
  \begin{example}
    \begin{itemize}
      \item $\sum_{k=1}^\infty \frac{x^k}{k}$ converges on $[-1, 1)$. Its derivative
        $\sum_{k=1}^\infty x^{k-1}$ converges on $(-1, 1)$ only.
      \item $\sum_{k=1}^\infty \frac{x^k}{k^2}$ converges on $[-1, 1]$. Its derivative
        $\sum_{k=1}^\infty \frac{x^{k-1}}{k}$ converges on $[-1, 1)$.
    \end{itemize}
  \end{example}
  Differentiating can lose an endpoint; integrating can gain one.
  The theorem is about the open interval $(a - R, a + R)$.
\end{frame}

% =============================================================
\section{New series from old}

\begin{frame}{Substitution into the geometric series}
  Start from $\frac{1}{1-u} = \sum_{k=0}^\infty u^k$ for $|u| < 1$:
  \begin{itemize}
    \item $u = -x$: \quad
      $\dfrac{1}{1+x} = \sum_{k=0}^\infty (-1)^k x^k$, \quad $|x| < 1$.
    \item $u = -x^2$: \quad
      $\dfrac{1}{1+x^2} = \sum_{k=0}^\infty (-1)^k x^{2k}$, \quad $|x| < 1$.
    \item $u = \frac{x}{2}$: \quad
      $\dfrac{1}{2-x} = \dfrac12 \cdot \dfrac{1}{1 - x/2} = \sum_{k=0}^\infty \dfrac{x^k}{2^{k+1}}$,
      \quad $|x| < 2$.
    \item Another centre, $u = -(x-1)$: \quad
      $\dfrac1x = \dfrac{1}{1 + (x-1)} = \sum_{k=0}^\infty (-1)^k (x-1)^k$, \quad $|x - 1| < 1$.
  \end{itemize}
  \pause
  {\small Note $\frac{1}{1+x^2}$ is defined for every $x$, yet its series has $R = 1$.
  Something at $x = \pm i$ in the complex plane decides that --- a story for complex analysis.\par}
\end{frame}

\begin{frame}{Integrating term by term}
  \begin{itemize}
    \item $\displaystyle\ln(1 + x) = \int_0^x \frac{\mathrm{d}t}{1+t}
      = \sum_{k=1}^\infty (-1)^{k+1} \frac{x^k}{k} = x - \frac{x^2}{2} + \frac{x^3}{3} - \cdots$
      \quad $(|x| < 1)$
    \item $\displaystyle -\ln(1 - x) = \sum_{k=1}^\infty \frac{x^k}{k}$ \quad $(|x| < 1)$
    \item $\displaystyle\arctan x = \int_0^x \frac{\mathrm{d}t}{1+t^2}
      = \sum_{k=0}^\infty (-1)^k \frac{x^{2k+1}}{2k+1} = x - \frac{x^3}{3} + \frac{x^5}{5} - \cdots$
      \quad $(|x| < 1)$
  \end{itemize}
  \pause
  \begin{remark}
    At $x = 1$ the series for $\ln(1+x)$ and $\arctan x$ converge (alternating series).
    By Abel's theorem their sums are $\ln 2$ and $\arctan 1 = \frac\pi4$. Next lecture
    proves both directly, with an estimate of the error.
  \end{remark}
\end{frame}

\begin{frame}{Differentiating term by term}
  From $\frac{1}{1-x} = \sum_{k=0}^\infty x^k$, for $|x| < 1$:
  \[ \frac{1}{(1-x)^2} = \sum_{k=1}^\infty k x^{k-1}, \qquad
     \frac{x}{(1-x)^2} = \sum_{k=1}^\infty k x^k, \qquad
     \frac{x(1+x)}{(1-x)^3} = \sum_{k=1}^\infty k^2 x^k . \]
  \pause
  \begin{example}[A binary counter]
    An increment flips the trailing $1$s and one $0$: $\texttt{1011} \to \texttt{1100}$ flips
    $3$ bits. A fraction $\frac{1}{2^k}$ of all values end in $\texttt{0}$ and then exactly $k - 1$
    ones, so the average number of bit flips per increment is
    $\sum_{k=1}^\infty \frac{k}{2^k} = \frac{1/2}{(1 - 1/2)^2} = 2$.
  \end{example}
\end{frame}

\begin{frame}{The exponential function}
  Let $E(x) = \sum_{k=0}^\infty \frac{x^k}{k!}$. Its radius is $\infty$, and term by term
  \[ E'(x) = \sum_{k=1}^\infty \frac{k\, x^{k-1}}{k!} = \sum_{k=1}^\infty \frac{x^{k-1}}{(k-1)!}
     = E(x), \qquad E(0) = 1 . \]
  \pause
  \begin{theorem}
    $e^x = \sum_{k=0}^\infty \frac{x^k}{k!}$ for every $x \in \R$.
  \end{theorem}
  {\small Proof: $\bigl(E(x) e^{-x}\bigr)' = E(x) e^{-x} - E(x) e^{-x} = 0$, so
  $E(x) e^{-x} = E(0) e^0 = 1$.\par}
  \pause
  \begin{itemize}
    \item $x = 1$: \ $e = \sum_{k=0}^\infty \frac{1}{k!}$, so $\sum_{k=1}^\infty \frac{1}{k!} = e - 1$
      (Lecture 2).
    \item $x = -1$: \ $\frac1e = \sum_{k=0}^\infty \frac{(-1)^k}{k!}$, so
      $\sum_{k=1}^\infty \frac{(-1)^{k+1}}{k!} = 1 - \frac1e$ (Lecture 3).
  \end{itemize}
\end{frame}

\begin{frame}{Multiplying power series}
  \begin{theorem}[Cauchy product]
    If $\sum a_k x^k$ and $\sum b_k x^k$ both converge for $|x| < R$, then for $|x| < R$
    \[ \Bigl(\sum_{k=0}^\infty a_k x^k\Bigr) \Bigl(\sum_{k=0}^\infty b_k x^k\Bigr)
       = \sum_{k=0}^\infty c_k x^k, \qquad c_k = \sum_{j=0}^{k} a_j b_{k-j} . \]
  \end{theorem}
  {\small Both series converge absolutely, so all products $a_j x^j \cdot b_m x^m$ may be added
  in any order (Lecture 3) --- in particular grouped by $j + m = k$. We omit the details.\par}
\end{frame}

\begin{frame}{Multiplying power series: an example}
  \begin{example}
    $\dfrac{1}{(1-x)^2} = \Bigl(\sum x^k\Bigr)^2 = \sum_{k=0}^\infty (k+1) x^k$, \ since
    $c_k = \sum_{j=0}^k 1 \cdot 1 = k + 1$. The same as by differentiation.
  \end{example}
  {\small The sequence $(c_k)$ is the \emph{convolution} of $(a_k)$ and $(b_k)$.
  In Unit 7 this turns counting problems into algebra.\par}
\end{frame}

% =============================================================
\section{Summary and practice}

\begin{frame}{Summary}
  \begin{itemize}
    \item A power series $\sum c_k (x - a)^k$ has a \textbf{radius of convergence}
      $R \in [0, \infty]$: absolute convergence for $|x - a| < R$, divergence for
      $|x - a| > R$. Check the endpoints separately.
    \item $R = 1/L$ if $\bigl|\frac{c_{k+1}}{c_k}\bigr| \to L$ or $\sqrt[k]{|c_k|} \to L$.
    \item Uniform convergence on every $[a - r, a + r]$ with $r < R$.
    \item Inside, differentiate and integrate term by term; the radius does not change.
    \item $c_k = \frac{f^{(k)}(a)}{k!}$: a function has at most one power series around $a$.
    \item New series from old: substitute, integrate, differentiate, multiply.
      Results so far: $\frac{1}{1-x}$, $\ln(1+x)$, $\arctan x$, $e^x$.
  \end{itemize}
\end{frame}

\begin{frame}{Practice}
  \small
  \begin{enumerate}
    \item Find the radius and the interval of convergence:
      \begin{enumerate}[(a)]
        \item $\sum_{k=1}^\infty \frac{x^k}{k\, 2^k}$ \qquad
          (b)\ $\sum_{k=1}^\infty k^2 x^k$ \qquad
          (c)\ $\sum_{k=1}^\infty \frac{(x+1)^k}{k^2}$ \qquad
          (d)\ $\sum_{k=0}^\infty \frac{(-1)^k x^{2k}}{(2k)!}$
      \end{enumerate}
    \item Find a power series for $\frac{x}{1+x^2}$ around $0$, and its radius.
    \item Find a power series for $\ln(2 + x)$ around $0$, and its radius.
    \item Find $\sum_{k=1}^\infty \frac{k}{3^k}$.
    \item Use $\arctan\frac{1}{\sqrt3} = \frac\pi6$ to write $\pi$ as a series. How fast does it
      converge?
    \item Show that if $\sum c_k x^k$ has radius $R$, then $\sum c_k x^{2k}$ has radius $\sqrt R$.
    \item Find $\bigl(\sum_{k=0}^\infty x^k\bigr)\bigl(\sum_{k=0}^\infty (-1)^k x^k\bigr)$
      as a Cauchy product, and check your answer another way.
  \end{enumerate}
\end{frame}

\begin{frame}{Answers}
  \small
  \begin{enumerate}
    \item (a) $R = 2$, interval $[-2, 2)$. \
      (b) $R = 1$, interval $(-1, 1)$. \
      (c) $R = 1$, interval $[-2, 0]$. \
      (d) $R = \infty$ (ratio of terms $\frac{x^2}{(2k+1)(2k+2)} \to 0$).
    \item $x \sum_{k=0}^\infty (-1)^k x^{2k} = \sum_{k=0}^\infty (-1)^k x^{2k+1}$, \ $R = 1$.
    \item $\ln 2 + \ln\bigl(1 + \frac x2\bigr) = \ln 2 + \sum_{k=1}^\infty (-1)^{k+1} \frac{x^k}{k\, 2^k}$,
      \ $R = 2$.
    \item $\frac{x}{(1-x)^2}$ at $x = \frac13$: \ $\frac{1/3}{4/9} = \frac34$.
    \item $\pi = 2\sqrt3 \sum_{k=0}^\infty \frac{(-1)^k}{3^k (2k+1)}$. The terms shrink by a
      factor of about $3$ each time: roughly one decimal digit for every two terms.
    \item $\sum c_k (x^2)^k$ converges absolutely if $x^2 < R$ and diverges if $x^2 > R$.
    \item $c_k = \sum_{j=0}^k (-1)^{k-j}$ is $1$ for even $k$ and $0$ for odd $k$: the product
      is $\sum x^{2k} = \frac{1}{1-x^2} = \frac{1}{1-x} \cdot \frac{1}{1+x}$.
  \end{enumerate}
\end{frame}

\begin{frame}{Next time}
  \textbf{Lecture 06 --- Taylor series} (Unit 6)
  \begin{itemize}
    \item Which functions are power series? Taylor and Maclaurin series
    \item Taylor's theorem and estimates of the remainder
    \item Series for $\sin$, $\cos$, $(1 + x)^\alpha$; and, at last, $\ln 2$ and $\frac\pi4$
  \end{itemize}
  \vskip1em
  {\small
  Reading: Strang \& Herman, \emph{Calculus, Volume 2}, Sections 6.3 and 6.4.\\[0.4em]
  Questions: the module channel on Discord, or office hours,
  Monday--Friday, 10:00--17:00 (Baku time, UTC+4),
  \href{mailto:ali.khudiyev@wosm.academy}{ali.khudiyev@wosm.academy}.\par}
\end{frame}

\end{document}
